\section*{Hoofdstuk \ref{ch:b2:enolates-aldol} --- Enolen, enolaten en de aldolreactie}

\begin{solution}{exo:b2:enolates-aldol:1}
Butanon: but-1-een-2-ol \ce{CH2=C(OH)CH2CH3} en but-2-een-2-ol \ce{CH3C(OH)=CHCH3} ((\textit{E}) en
(\textit{Z})). Cyclohexanon: cyclohex-1-een-1-ol.
\end{solution}

\begin{solution}{exo:b2:enolates-aldol:2}
Ethaan $<$ propanon $<$ ethyl-3-oxobutanoaat (10.7) $<$ pentaan-2,4-dion (9.0).
\end{solution}

\begin{solution}{exo:b2:enolates-aldol:3}
3-Hydroxybutanal, en bij verhitten daarna but-2-enal.
\end{solution}

\begin{solution}{exo:b2:enolates-aldol:4}
Diethylpropaandioaat, natriumethoxide, 1-broombutaan; hydrolyse tot butylpropaandizuur;
verhitten verwijdert \ce{CO2}: \ce{CH3(CH2)3CH2COOH}, hexaanzuur.
\end{solution}

\begin{solution}{exo:b2:enolates-aldol:5}
Kinetisch: deprotonering op C6 (de kant van \ce{CH2}), met LDA bij \qty{-78}{\degreeCelsius}. Thermodynamisch:
de meer gesubstitueerde dubbele binding C1=C2, naar het koolstofatoom met de methylgroep, met een alkoxide in zijn
alcohol bij kamertemperatuur.
\end{solution}

\begin{solution}{exo:b2:enolates-aldol:6}
Het product is 4-fenylbut-3-een-2-on (benzalaceton), \ce{C6H5CH=CHCOCH3}. Benzaldehyde heeft geen $\alpha$-waterstofatoom: het
kan geen \omterm{def:b2:enolates-aldol:enolate}{enolaat} vormen en is alleen het elektrofiel.
\end{solution}

\begin{solution}{exo:b2:enolates-aldol:7}
Het \omterm{def:b2:enolates-aldol:enolate}{enolaat} van een van de methylgroepen (C1) valt de andere carbonylgroep (C5) aan: een vijfring,
na dehydratatie 3-methylcyclopent-2-een-1-on. Het alternatief, het \omterm{def:b2:enolates-aldol:enolate}{enolaat} op C3 dat C5 aanvalt,
zou een gespannen driering geven.
\end{solution}

\begin{solution}{exo:b2:enolates-aldol:8}
\ce{2CH3COOC2H5 -> CH3COCH2COOC2H5 + C2H5OH} met natriumethoxide; het gedeprotoneerde product wordt
bij aanzuren teruggewonnen als ethyl-3-oxobutanoaat.
\end{solution}

\begin{solution}{exo:b2:enolates-aldol:9}
Propanon, ethanal en acetofenon (die hebben een groep \ce{CH3-CO}; ethanal is \ce{CH3CHO}).
Pentaan-3-on niet.
\end{solution}

\begin{solution}{exo:b2:enolates-aldol:10}
Het \omterm{def:b2:enolates-aldol:enolate}{enolaat} aan het ene uiteinde valt de ester aan het andere aan: ethyl-2-oxocyclopentaan-1-carboxylaat, een
vijfring.
\end{solution}

\begin{solution}{exo:b2:enolates-aldol:11}
In 3-methylpentaan-2-on is het stereocentrum (C3) het $\alpha$-koolstofatoom: enolisering maakt het vlak, en
herprotonering vanaf een van beide vlakken racemiseert het. In 4-methylhexaan-2-on zit het stereocentrum op C4,
$\beta$ ten opzichte van de carbonylgroep, en blijft het bij enolisering onaangeroerd.
\end{solution}

\begin{solution}{exo:b2:enolates-aldol:12}
Propanon met twee equivalenten benzaldehyde in base. Met twee verschillende aldehyden condenseer je
propanon eerst met het ene aldehyde (één equivalent, en benzalaceton isoleren), en daarna zijn
resterende methylgroep met het tweede. In één pot concurreren de twee aldehyden en ontstaat een mengsel van
drie dienonen.
\end{solution}

\begin{solution}{pb:b2:enolates-aldol:1}
\textbf{1.} Propanon is de enige partner met $\alpha$-waterstofatomen: zes, drie op elke methylgroep.
\textbf{2.} $\mathrm{CH_3{-}CO{-}CH_2^-} \leftrightarrow \mathrm{CH_3{-}C(O^-){=}CH_2}$.
\textbf{3.} Het heeft geen waterstofatoom op het koolstofatoom naast zijn carbonylgroep (dat koolstofatoom hoort bij de ring en
draagt geen H).
\textbf{4.} Benzaldehyde, een aldehyde in overmaat, is het betere elektrofiel; het \omterm{def:b2:enolates-aldol:enolate}{enolaat} van propanon ontmoet het
vaker dan een carbonylgroep van propanon.
\textbf{5.} Een keton heeft twee koolstofgroepen die elektronendichtheid naar zijn carbonylkoolstofatoom duwen en
de nadering hinderen; een aldehyde heeft er één.
\textbf{6.} Er is maar een kleine fractie \omterm{def:b2:enolates-aldol:enolate}{enolaat} nodig: het wordt teruggevormd naarmate het reageert.
\textbf{7.} Het enolaatkoolstofatoom bindt aan het carbonylkoolstofatoom van benzaldehyde: het alkoxide van
4-hydroxy-4-fenylbutaan-2-on.
\textbf{8.} Wegname van een $\alpha$-waterstofatoom en verlies van hydroxide van het $\beta$-koolstofatoom: 4-fenylbut-3-een-2-on.
\textbf{9.} De gevormde dubbele binding is geconjugeerd met zowel de ring als de carbonylgroep.
\textbf{10.} De andere methylgroep heeft nog drie $\alpha$-waterstofatomen.
\textbf{11.} \ce{2C6H5CHO + CH3COCH3 -> C6H5CH=CHCOCH=CHC6H5 + 2H2O}.
\textbf{12.} De dehydrataties, en het neerslaan van het product, dat de oplossing verlaat.
\textbf{13.} Twee, allebei (\textit{E}).
\textbf{14.} De (\textit{Z})-isomeren zetten de fenylgroep naast de carbonylgroep: sterische spanning.
\textbf{15.} Twee \ce{C=C}, één \ce{C=O} en de twee ringen: een geconjugeerd systeem over het hele molecuul.
\textbf{16.} Zijn uitgebreide conjugatie verkleint de $\pi$-kloof (\cref{prop:b2:huckel:gap}): het absorbeert aan het
blauwe uiteinde van het zichtbare en ziet er geel uit.
\textbf{17.} Nee: het is vlak (of bijna) met een spiegelvlak.
\textbf{18.} 106.12, 58.08 en \qty{234.30}{g/mol}.
\textbf{19.} Benzaldehyde $2.1 \times 1.045/106.12 = \qty{20.7}{mmol}$; propanon $0.75 \times
0.7845/58.08 = \qty{10.1}{mmol}$.
\textbf{20.} Propanon vraagt \qty{20.3}{mmol} benzaldehyde, en er is \qty{20.7}{mmol} beschikbaar:
propanon is beperkend, maar net.
\textbf{21.} $0.01013 \times 234.30 = \qty{2.37}{g}$.
\textbf{22.} Overgebleven benzaldehyde blijft in de ethanol en wordt weggewassen.
\textbf{23.} Het monocondensatieproduct, benzalaceton.
\textbf{24.} $0.75 \times 2.37 = \boldsymbol{\approx \qty{1.78}{g}}$ dibenzalaceton.
\end{solution}
